% Part: first-order-logic
% Chapter: completeness
% Section: construction-of-model

\documentclass[../../../include/open-logic-section]{subfiles}

\begin{document}

\iftag{FOL}
      {\olfileid[fr]{fol}{com}{mod}}
      {\olfileid[fr]{pl}{com}{mod}}

\olsection{Construction d'un modèle}

\begin{explain}
\iftag{FOL}{Nous laissons pour l'instant de côté~$\eq$ : nous voulons
  seulement montrer qu'un ensemble cohérent~$\Gamma$ d'!!{sentence}s
  où $\eq$ ne figure pas est satisfaisable. Étendons d'abord
  $\Gamma$ en un ensemble cohérent, !!{complete} et saturé
  $\Gamma^*$. La définition du modèle~$\Struct{M(\Gamma^*)}$ est
  alors simple. Prenons pour domaine l'ensemble des termes clos de
  $\Lang{L'}$. Interprétons chaque !!{constant} par elle-même, et
  plus généralement chaque terme clos~$t$ de sorte que
  $\Value{t}{M(\Gamma^*)}=t$. Attribuons aux !!{predicate}s des
  extensions telles qu'un !!{sentence} atomique soit vrai dans
  $\Struct{M(\Gamma^*)}$ si et seulement s'il appartient à
  $\Gamma^*$. Ainsi, les !!{sentence}s atomiques de~$\Gamma^*$
  seront vrais dans ce modèle. Les autres le seront grâce à la
  cohérence, à la complétude et à la saturation de~$\Gamma^*$.}{Nous
  pouvons maintenant définir !!a{valuation} qui rend vraies toutes
  les !!{formula}s~$!A\in\Gamma$. Appliquons d'abord le lemme de
  Lindenbaum pour obtenir un ensemble cohérent et complet
  $\Gamma^*\supseteq\Gamma$. Les !!{propositional variable}s de
  $\Gamma^*$ détermineront la valuation~$\pAssign v(\Gamma^*)$.}
\end{explain}

\iftag{FOL}{%
\begin{defn}[Modèle des termes]
\ollabel{defn:termmodel}
Soit $\Gamma^*$ un ensemble cohérent, !!{complete} et saturé d'%
!!{sentence}s dans un langage~$\Lang L$ contenant au moins une
constante, comme après l'expansion de Henkin. Le \emph{modèle des
termes}~$\Struct M(\Gamma^*)$ associé à~$\Gamma^*$ est la
!!{structure} définie ainsi.
\begin{enumerate}
\item Le !!{domain}~$\Domain{M(\Gamma^*)}$ est l'ensemble de tous les
  termes clos de~$\Lang L$.
\item Une !!{constant}~$c$ s'interprète par elle-même :
  $\Assign{c}{M(\Gamma^*)}=c$.
\item À un !!{function}~$f$ d'arité~$n$, on attribue la
  fonction qui, aux termes clos $t_1,\dots,t_n$, associe le terme
  clos $f(t_1,\dots,t_n)$ :
\[
\Assign{f}{M(\Gamma^*)}(t_1, \dots, t_n) = f(t_1,\dots, t_n)
\]
\item Si $R$ est un !!{predicate} d'arité~$n$, alors
  \[
  \tuple{t_1, \dots,
  t_n} \in \Assign{R}{M(\Gamma^*)} \quad\Longleftrightarrow\quad \Atom{R}{t_1, \dots,
    t_n} \in \Gamma^*.
  \]
\end{enumerate}
\end{defn}
}{%
\begin{defn}
\ollabel{defn:termmodel}
Soit $\Gamma^*$ un ensemble cohérent et complet de !!{formula}s.
Nous définissons la valuation
\[
\pAssign v(\Gamma^*)(p) = \begin{cases}
  \True & \text{si $p \in \Gamma^*$}\\
  \False & \text{si $p \notin \Gamma^*$}
  \end{cases}
\]
\end{defn}
}

\iftag{FOL}{%
Vérifions que $\Value{t}{M(\Gamma^*)}=t$ pour tout terme clos~$t$.

\begin{lem}
\ollabel{lem:val-in-termmodel} Soit $\Struct M(\Gamma^*)$ le modèle
des termes de la \olref{defn:termmodel}. Alors
$\Value{t}{M(\Gamma^*)}=t$ pour tout terme clos~$t$.
\end{lem}
\begin{proof}
Raisonnons par récurrence sur la construction de~$t$. Si $t$ est une
!!{constant}, l'égalité découle directement de la définition du
modèle. Supposons ensuite que $t_1,\ldots,t_n$ soient des termes clos
vérifiant $\Value{t_i}{M(\Gamma^*)}=t_i$, et que $f$ soit un
!!{function} d'arité~$n$. Alors
\begin{align*}
\Value{f(t_1,\ldots,t_n)}{M(\Gamma^*)} &= \Assign{f}{M(\Gamma^*)}(\Value{t_1}
 {M(\Gamma^*)},
\ldots, \Value{t_n}{M(\Gamma^*)}) \\
&= \Assign{f}{M(\Gamma^*)}(t_1, \dots, t_n) \\
&= f(t_1,\dots, t_n).
\end{align*}
L'égalité vaut donc pour tout terme clos~$t$.
\end{proof}
}

\iftag{FOL}{%
\begin{explain}
  \iftag{prvEx}{Une !!{structure}~$\Struct{M}$ peut rendre vrai un
    !!{sentence} existentiel~$\lexists[x][!A(x)]$ sans rendre
    vraie aucune instance~$!A(t)$.}{}
  \iftag{prvAll}{Une !!{structure}~$\Struct{M}$ peut rendre vraies
    toutes les instances~$!A(t)$ d'un !!{sentence} universel
    $\lforall[x][!A(x)]$ sans rendre vrai ce dernier.}{}
  En général, tout !!{element} du !!{domain}~$\Domain{M}$ n'est pas
  nécessairement la valeur d'un terme clos : la !!{structure} peut
  ne pas être couverte par les termes. C'est pourquoi la relation de
  satisfaction est définie à l'aide d'assignations aux variables.
  Dans le modèle des termes~$\Struct{M(\Gamma^*)}$, cependant, chaque
  !!{element} est désigné par un terme clos. Le résultat suivant
  précise cette propriété.
\end{explain}

\begin{prop}
\ollabel{prop:quant-termmodel}
Soit $\Struct M(\Gamma^*)$ le modèle des termes de la
\olref{defn:termmodel}. Supposons que les !!{variable}s libres
de~$!A(x)$ appartiennent à~$\{x\}$.
\begin{tagenumerate}{prvEx,prvAll}
\tagitem{prvEx}{$\Sat{M(\Gamma^*)}{\lexists[x][!A(x)]}$ si et
  seulement si $\Sat{M(\Gamma^*)}{!A(t)}$ pour au moins un terme
  clos~$t$.}{}
\tagitem{prvAll}{$\Sat{M(\Gamma^*)}{\lforall[x][!A(x)]}$ si et
  seulement si $\Sat{M(\Gamma^*)}{!A(t)}$ pour tout terme clos~$t$.}{}
\end{tagenumerate}
\end{prop}

\begin{proof}
\begin{tagenumerate}{prvEx,prvAll}
\tagitem{prvEx}{%
  \iftag{probEx}{Exercice.}{D'après la
    \olref[syn][ass]{prop:sat-quant},
    $\Sat{M(\Gamma^*)}{\lexists[x][!A(x)]}$ si et seulement s'il
    existe une assignation~$s$ telle que
    $\Sat{M(\Gamma^*)}{!A(x)}[s]$. Comme
    $\Domain{M(\Gamma^*)}$ est l'ensemble des termes clos de
    $\Lang{L}$, cela revient à dire qu'il existe une assignation~$s$
    et un terme clos~$t$ tels que $s(x)=t$ et
    $\Sat{M(\Gamma^*)}{!A(x)}[s]$. La valeur de~$t$ est~$t$.
    Par la \olref[fol][syn][ext]{prop:ext-formulas}, cette dernière
    assertion équivaut à $\Sat{M(\Gamma^*)}{!A(t)}[s]$ lorsque
    $s(x)=t$. Enfin, la
    \olref[fol][syn][ass]{prop:sentence-sat-true} donne
    $\Sat{M(\Gamma^*)}{!A(t)}[s]$ si et seulement si
    $\Sat{M(\Gamma^*)}{!A(t)}$, puisque $!A(t)$ est un
    !!{sentence}.}}{}
\tagitem{prvAll}{%
  \iftag{probAll}{Exercice.}{D'après la
    \olref[syn][ass]{prop:sat-quant},
    $\Sat{M(\Gamma^*)}{\lforall[x][!A(x)]}$ si et seulement si,
    pour toute assignation~$s$,
    $\Sat{M(\Gamma^*)}{!A(x)}[s]$. Le domaine étant constitué des
    termes clos de~$\Lang{L}$, chaque valeur~$s(x)$ est un tel
    terme~$t$, et tout terme clos peut être choisi comme valeur de
    $x$. La valeur du terme~$t$ est~$t$. Par la
    \olref[fol][syn][ext]{prop:ext-formulas},
    $\Sat{M(\Gamma^*)}{!A(x)}[s]$ si et seulement si
    $\Sat{M(\Gamma^*)}{!A(t)}[s]$ pour $s(x)=t$. Par la
    \olref[fol][syn][ass]{prop:sentence-sat-true}, cette dernière
    condition équivaut à $\Sat{M(\Gamma^*)}{!A(t)}$, car $!A(t)$
    est un !!{sentence}.}}{}
\end{tagenumerate}
\end{proof}
}{}

\tagprob[FOL]{probEx,probAll}
\begin{prob}
Achever la démonstration de la
\olref[fol][com][mod]{prop:quant-termmodel}.
\end{prob}
\tagendprob

\begin{lem}[Lemme de vérité]
\ollabel{lem:truth} Soit $!A$ \iftag{FOL}{un !!{sentence}}{une !!{formula}}.
\iftag{FOL}{Supposons que $!A$ ne contienne pas
$\eq$. Alors}{}
$\iftag{FOL}{\Sat{M(\Gamma^*)}}{\pSat{v(\Gamma^*)}}{!A}$ si et
seulement si $!A\in\Gamma^*$.
\end{lem}

\begin{proof}
Démontrons les deux sens simultanément par récurrence sur le nombre
de connecteurs\iftag{FOL}{ et de quantificateurs}{} de~$!A$.
\iftag{FOL}{La substitution d'un terme clos dans une formule conserve
ce nombre. Dans un cas quantifié, les instances de la matrice ont donc
un opérateur logique de moins que l'énoncé quantifié.}{}
\begin{enumerate}
\tagitem{prvFalse}{\indcase{!A}{\lfalse}
  {$\iftag{FOL}{\Sat/{M(\Gamma^*)}{\lfalse}}{\pSat/{v(\Gamma^*)}{\lfalse}}$
    par définition de la satisfaction. D'autre part,
    $\lfalse\notin\Gamma^*$ puisque $\Gamma^*$ est cohérent.}}{}

\tagitem{prvTrue}{\indcase{!A}{\ltrue}
  {$\iftag{FOL}{\Sat{M(\Gamma^*)}}{\pSat{v(\Gamma^*)}}{\ltrue}$
    par définition de la satisfaction. D'autre part,
    $\ltrue\in\Gamma^*$ car $\Gamma^*$ est cohérent et
    !!{complete}, et $\Gamma^*\Proves\ltrue$.}}{}

\tagitem{FOL}{\indcase{!A}{R(t_1, \dots, t_n)}
  {$\Sat{M(\Gamma^*)}{\Atom{R}{t_1, \dots, t_n}}$ si et seulement si
    $\tuple{t_1,\dots,t_n}\in\Assign{R}{M(\Gamma^*)}$, par
    définition de la satisfaction et par le
    \olref{lem:val-in-termmodel}, si et seulement si
    $R(t_1,\dots,t_n)\in\Gamma^*$, par la construction de
    $\Struct M(\Gamma^*)$.}}{\indcase{!A}{p}
  {$\pSat{v(\Gamma^*)}{p}$ si et seulement si
    $\pAssign v(\Gamma^*)(p)=\True$, par définition de la
    satisfaction, si et seulement si $p\in\Gamma^*$, par la
    construction de~$\pAssign v(\Gamma^*)$.}}

\tagitem{prvNot}{%
  \indcase{!A}{\lnot !B}
    {$\iftag{FOL}{\Sat{M(\Gamma^*)}}{\pSat{v(\Gamma^*)}}{\indfrm}$
    si et seulement si
    $\iftag{FOL}{\Sat/{M(\Gamma^*)}{!B}}{\pSat/{v(\Gamma^*)}{!B}}$,
    par définition de la satisfaction. Par hypothèse de récurrence,
    cette dernière condition équivaut à $!B\notin\Gamma^*$.
    Puisque $\Gamma^*$ est cohérent et !!{complete},
    $!B\notin\Gamma^*$ si et seulement si
    $\lnot !B\in\Gamma^*$.}}{}

\tagitem{prvAnd}{%
\iftag{probAnd}{%
  \indcase!{!A}{!B \land !C}{}}{%
  \indcase{!A}{!B \land !C}
    {$\iftag{FOL}{\Sat{M(\Gamma^*)}}{\pSat{v(\Gamma^*)}}{\indfrm}$
    si et seulement si
    $\iftag{FOL}{\Sat{M(\Gamma^*)}}{\pSat{v(\Gamma^*)}}{!B}$ et
    $\iftag{FOL}{\Sat{M(\Gamma^*)}}{\pSat{v(\Gamma^*)}}{!C}$,
    par définition de la satisfaction, si et seulement si
    $!B\in\Gamma^*$ et $!C\in\Gamma^*$, par hypothèse de
    récurrence. D'après la
    \olref[ccs]{prop:ccs}\olref[ccs]{prop:ccs-and}, cela équivaut à
    $(!B\land !C)\in\Gamma^*$.}}}{}

\tagitem{prvOr}{%
\iftag{probOr}{%
  \indcase!{!A}{!B \lor !C}{}}{%
  \indcase{!A}{!B \lor !C}
    {$\iftag{FOL}{\Sat{M(\Gamma^*)}}{\pSat{v(\Gamma^*)}}{\indfrm}$
    si et seulement si
    $\iftag{FOL}{\Sat{M(\Gamma^*)}}{\pSat{v(\Gamma^*)}}{!B}$ ou
    $\iftag{FOL}{\Sat{M(\Gamma^*)}}{\pSat{v(\Gamma^*)}}{!C}$,
    par définition de la satisfaction, si et seulement si
    $!B\in\Gamma^*$ ou $!C\in\Gamma^*$, par hypothèse de
    récurrence. D'après la
    \olref[ccs]{prop:ccs}\olref[ccs]{prop:ccs-or}, cela équivaut à
    $(!B\lor !C)\in\Gamma^*$.}}}{}

\tagitem{prvIf}{%
\iftag{probIf}{%
  \indcase!{!A}{!B \lif !C}{}}{%
  \indcase{!A}{!B \lif !C}
    {$\iftag{FOL}{\Sat{M(\Gamma^*)}}{\pSat{v(\Gamma^*)}}{\indfrm}$
    si et seulement si
    $\iftag{FOL}{\Sat/{M(\Gamma^*)}}{\pSat/{v(\Gamma^*)}}{!B}$ ou
    $\iftag{FOL}{\Sat{M(\Gamma^*)}}{\pSat{v(\Gamma^*)}}{!C}$,
    par définition de la satisfaction, si et seulement si
    $!B\notin\Gamma^*$ ou $!C\in\Gamma^*$, par hypothèse de
    récurrence. D'après la
    \olref[ccs]{prop:ccs}\olref[ccs]{prop:ccs-if}, cela équivaut à
    $(!B\lif !C)\in\Gamma^*$.}}}{}

\iftag{FOL}{%
\tagitem{prvAll}{%
  \iftag{probAll}{%
    \indcase!{!A}{\lforall[x][!B(x)]}{}}{%
    \indcase{!A}{\lforall[x][!B(x)]}{%
      $\Sat{M(\Gamma^*)}{\indfrm}$ si et seulement si
      $\Sat{M(\Gamma^*)}{!B(t)}$ pour tout terme clos~$t$,
      d'après la \olref{prop:quant-termmodel}. Par hypothèse de
      récurrence, cela équivaut à $!B(t)\in\Gamma^*$ pour tout
      terme clos~$t$. Par la
      \olref[hen]{prop:saturated-instances}, cela équivaut à
      $\lforall[x][!B(x)]\in\Gamma^*$.}}}{}

\tagitem{prvEx}{%
  \iftag{probEx}{%
    \indcase!{!A}{\lexists[x][!B(x)]}{}}{%
    \indcase{!A}{\lexists[x][!B(x)]}{%
      $\Sat{M(\Gamma^*)}{\indfrm}$ si et seulement si
      $\Sat{M(\Gamma^*)}{!B(t)}$ pour au moins un terme clos~$t$,
      d'après la \olref{prop:quant-termmodel}. Par hypothèse de
      récurrence, cela équivaut à $!B(t)\in\Gamma^*$ pour au
      moins un terme clos~$t$. Par la
      \olref[hen]{prop:saturated-instances}, cela équivaut à
      $\lexists[x][!B(x)]\in\Gamma^*$.}}}{}
}{}
\end{enumerate}
\end{proof}

\tagprob[FOL]{probOr,probAnd,probIf,probEx,probAll}

\begin{prob}
Achever la démonstration du \olref[fol][com][mod]{lem:truth}.
\end{prob}

\tagendprob

\tagprob[notFOL]{probOr,probAnd,probIf}

\begin{prob}
Achever la démonstration du \olref[pl][com][mod]{lem:truth}.
\end{prob}

\tagendprob

\end{document}
